
\chapter{Сверка}\label{ch:reconciliation}

\highlight{Когда деньги \enquote{теряются}, они обычно расходятся между несколькими записями}. В~журнале авторизаций покупка уже разрешена, в~клиринговом файле она ещё не появилась, в~банковской выписке от неё осталась только нетто-сумма после комиссий и~взаимных вычетов. Сверка собирает из этих разрозненных следов одну операцию.

Внутренний реестр обязательств из главы~\ref{ch:internal-ledger}
отвечает на вопрос, какое денежное состояние видит сам продукт.
Сверка отвечает на следующий вопрос, совпадает ли эта внутренняя картина
с~внешними файлами схемы, эквайера и~банка. Если перепутать эти два уровня, команда начнёт сравнивать несопоставимые величины и~искать
потери там, где есть разница между моделями учёта.

\section{Три разрыва}\label{sec:recon-gaps}

\highlight{Карточная операция образует цепочку событий, и~каждая стадия
оставляет собственный след. Авторизация фиксирует право провести операцию,
клиринг закрепляет её окончательный денежный облик, а~расчёт показывает
уже фактическое движение средств} с~комиссиями и взаимными вычетами.
Сверка начинается с~признания того, что эти следы по определению не обязаны
быть идентичными.

Первый разрыв возникает между авторизацией и клирингом. Ресторан может
авторизовать счёт на одну сумму, а~после чаевых отправить в~клиринг
другую; отель сначала блокирует предварительный депозит, а~потом
предъявляет итоговую стоимость проживания; авиакомпания доначисляет
часть сборов уже после исходной покупки. Поэтому сверка на этом уровне
начинается с~поиска правильной
пары. Для каждой клиринговой записи надо найти соответствующую
авторизацию и проверить, укладывается ли расхождение в~допустимые
пределы~\cite{visa-tadg,mc-tpr}.

Второй разрыв возникает между клирингом и расчётной позицией. За день
может пройти сто транзакций, но в~расчётном отчёте эквайер увидит
свёрнутую итоговую сумму. Здесь многие команды впервые чувствуют, что
\enquote{деньги не сходятся}, хотя сравнивают поток
операций с~результатом его нетто-свёртки. В~публичных правилах Mastercard нетто-расчёт описан как
базовая модель, а~прямой двусторонний расчёт между участниками вынесен
в~отдельное
исключение; Visa отдельно описывает VisaNet
Settlement Service как сервис, который формирует нетто-позицию расчёта
для участников, работающих через BASE~II (файловый клиринговый сервис
VisaNet, процессинговой сети Visa)~\cite{mc-rules,visa-core-rules,visa-visanet-booklet}.

Третий разрыв возникает между расчётной позицией и банковской выпиской.
Расчётная позиция в~системе схемы может уже считаться закрытой, а~движение по
корреспондентскому счёту пройти позже, в~другом окне дня или даже в
иной валюте после пересчёта. Здесь сверка идёт по сопоставлению
расчётного файла и~банковской выписки по счёту урегулирования;
отдельные транзакции на этом уровне уже не рассматриваются. Прежде чем искать
потери, назовите уровень, на котором их ищете. Иначе спор
превратится в~обсуждение разных способов описать один и~тот же поток.
Mastercard Rules требуют ежедневно сверять внутренние записи с
расчётными отчётами схемы и балансировать расхождения на этом уровне~\cite{mc-rules}.

Три разрыва за неделю — три параллельные истории.

История~A (первый разрыв). День~1: ресторан авторизует 5\,000~руб.,
холд создан. День~2: клиринговый файл приходит с~суммой 5\,200~руб.
(чаевые 200~руб.); сверка сопоставляет запись по ARN (Acquirer
Reference Number) и~коду авторизации, но фиксирует расхождение в 200~руб.

История~B (\textit{orphan authorization} -- авторизация без последующего
клиринга). День~1: авторизация на 800~руб. в~другом магазине. День~5:
клиринг на 800~руб. так и~не пришёл; срок представления (\textit{presentment
window} -- период, в~течение которого клиринговая запись должна быть
передана сети после авторизации) ещё не истёк, запись висит в~очереди
ожидания. День~8: срок представления для MCC (Merchant Category Code --
четырёхзначный код категории торгово-сервисного предприятия) магазина
(7~дней) истекает; запись перемещается в~очередь исключений как
\textit{orphan authorization}. Холд снимается автоматически.

История~C (второй и третий разрывы). День~3: в~расчётном файле схемы
нетто-позиция за день составляет 4\,800~руб. (5\,200~руб. клиринга
минус 400~руб. встречных обязательств). Это второй разрыв. День~4:
банковская выписка показывает зачисление 4\,795~руб. Разница в~5~руб.
объясняется комиссией за расчётное обслуживание, которая не видна
в~файле схемы. Это третий разрыв.

С~марта 2022~года глобальный карточный контур Visa и~Mastercard для российских эмитентов закрыт: их карты обрабатываются только во~внутрироссийском контуре через ОПКЦ~НСПК~\cite{cbr-ps,nspk-mir-rules}. Разделы~\ref{sec:recon-base2-ipm}--\ref{sec:recon-dcc} разбирают сверочные форматы и~правила схем в~их международной редакции; разделы~\ref{sec:recon-mir}--\ref{sec:recon-sbp} --- сверку в~российском контуре через НСПК, кобейджи и~СБП. Для российских эквайеров, обслуживающих карты иностранных эмитентов, и~для дочерних структур российских банков за~рубежом обе части остаются практическими.

\section{BASE II и IPM}\label{sec:recon-base2-ipm}

Клиринговые файлы схем -- рабочие документы клирингового процесса,
а~не бухгалтерский отчёт. Читайте их как машинный след
завершённых транзакций. У Visa этот формат в~публичных правилах
описан через BASE~II, у Mastercard через IPM (Integrated Product
Messages -- семейство клиринговых форматов Mastercard, исторически
обслуживаемых Global Clearing Management System,
GCMS)~\cite{visa-core-rules,mc-pfmi-disclosure,mc-switching-services}.
Для российских банков обработка карт Visa и~Mastercard с~марта 2022~года ведётся через ОПКЦ~НСПК (раздел~\ref{sec:recon-mir}) с~собственным клирингом; BASE~II и~IPM остаются актуальными для эквайеров, обслуживающих карты иностранных эмитентов, и~для дочерних структур российских банков за~рубежом~\cite{cbr-ps,nspk-mir-rules}.

Visa определяет BASE II как файловый сервис VisaNet, который собирает,
проводит клиринг, рассчитывает и доставляет файлы финансовой и нефинансовой
активности между Visa и участниками. В~приложении к~публичным правилам
перечислены семейства форматов BASE~II: Transaction Code (TC) групп
01--48 и 50--92 (каждая запись делится на подзаписи Transaction Code
Record, TCR0--TCR7), а~отдельно опубликованы материалы по VisaNet
Settlement Service,
который выдаёт нетто-позицию расчёта и расчётные отчёты для участников
BASE II~\cite{visa-core-rules}.

Mastercard, в~отличие от Visa, публично раскрывает именно термин
Integrated Product Messages (IPM) Clearing Formats как часть своих
стандартов, а Transaction Processing Rules регулярно отсылают к~полям
сообщения First Presentment/1240, описываемым в~руководстве
IPM Clearing Formats~\cite{mc-pfmi-disclosure,mc-tpr}. Из этого следует
для сверки, что клиринг живёт не в~том же сообщении, что авторизация,
а~поля и~коды Visa BASE II не совпадают с~полями и~кодами Mastercard;
делать вид, что совпадают, нельзя.

Visa BASE II и Mastercard IPM сообщают по сути одно и~то же: факт
завершённой операции с~её денежным обликом. Но структура полей
различается. В~BASE~II идентификаторы сопоставления живут
в~Transaction Code (TC) и~связке ARN с~RRN (Retrieval Reference Number);
в~IPM/1240 аналогичную роль играют DE~31 (Data Element~31 -- Acquirer
Reference Data) и~DE~48 (Additional Data) сообщения ISO~8583/IPM
(см.~\cite{iso-8583-2023})~\cite{visa-core-rules,mc-tpr}.
Валюта транзакции и валюта расчёта кодируются в~разных полях, а
формат представления суммы (с~разделителем или без, с~фиксированной
точностью или плавающей) тоже отличается между схемами.
Семантически одинаковые поля сопоставляются напрямую: Account Number $\leftrightarrow$ DE~2 (PAN), Destination Amount $\leftrightarrow$ DE~4 (Amount, Transaction), Destination Currency Code $\leftrightarrow$ DE~49 (Currency, Transaction), Auth Code $\leftrightarrow$ DE~38 (Approval Code), ARN $\leftrightarrow$ DE~31 (Acquirer Reference Data), MCC $\leftrightarrow$ DE~26 (Card Acceptor Business Code). Внутренняя структура 23-значного ARN одинакова в~обеих схемах: позиция~1 как \textit{mixed use}, 2--7 -- BIN эквайера, 8--11 -- дата формата YDDD (год + порядковый день года), 12--22 -- порядковый номер, 23 -- проверка по алгоритму Луна. Конкретный состав DE для IPM~1240 First Presentment и~подтверждённые позиции TCR~0 для Visa TC~05 в~TC~33.A показаны на~рис.~\ref{fig:base2-vs-ipm}; источники~--- Visa Acceptance Developer Portal~\cite{visa-acceptance-capture-spec} и~Mastercard GCMS Reference Manual~2021~\cite{mc-gcms-rm-2021}.

Прежде чем механизм сверки начнёт сопоставлять записи, приведите каждую
клиринговую запись к~единой внутренней модели: общий идентификатор операции,
сумма в~минорных единицах валюты, дата обработки в~одном часовом
поясе, код валюты по ISO~4217 (международный стандарт трёхбуквенных
и~трёхзначных кодов валют)~\cite{iso-4217}. Без этого шага даже идеальный алгоритм
сопоставления будет регулярно спотыкаться на разнице форматов;
реальные расхождения уйдут на второй план.

\begin{figure}[tbp]
  \centering
  \includefiguresvg[width=.8\linewidth]{assets/figures/ch30-reconciliation/base2-vs-ipm}
  \caption{Visa TC~05 в~TC~33.A (TCR~0) и~Mastercard IPM~1240 First Presentment: состав полей.}\label{fig:base2-vs-ipm}
\end{figure}

Парсер должен собрать клиринговую запись, извлечь поля сопоставления
и передать их в~механизм сверки. Сводить Visa BASE II и Mastercard IPM
к~одному формату нельзя. Для каждой схемы нужен свой адаптер или
хотя бы аккуратное промежуточное представление. Иначе ошибка возникает ещё до
бухгалтерии, на этапе разбора файла.

\section{Нетто- и~валовой расчёт}\label{sec:recon-net-vs-gross}

Сумма транзакций редко совпадает с~тем, что пришло на
расчётный счёт, потому что карточные системы показывают результат
нетто-свёртки; вала отдельных операций в~ней нет. По умолчанию работает
нетто-расчёт, валовой остаётся исключением.

\highlight{В~модели нетто-расчёта (\enquote{netting} -- итоговое
сальдо взаимных обязательств за период) схема сводит все позиции дня
в~одну итоговую сумму, и~за множеством покупок, возвратов и~комиссий
остаётся единое сальдо нетто-расчёта, которое и~видит эквайер}. Подход
уменьшает количество отдельных переводов и~делает расчётную
инфраструктуру дешевле и~проще в~эксплуатации.

Валовой расчёт (\enquote{gross settlement} -- расчёт каждой операции
по полной сумме отдельно, без свёртки встречных обязательств) устроен
иначе: каждая операция или небольшая группа операций идёт отдельно. Режим прозрачнее, потому что деньги не
растворяются в~общей дневной сумме, но платит он другую цену:
учёт становится более детальным, а дисциплина сверки более строгой.

Выбор модели расчёта не принадлежит ТСП, так как базовую
модель для участника определяет схема, а~эквайер транслирует её
в~условия договора. Mastercard описывает нетто-расчёт как стандартный; прямой
двусторонний расчёт возможен, но требует отдельного соглашения между
участниками~\cite{mc-rules}. Visa формирует нетто-позицию через
VisaNet Settlement Service~\cite{visa-core-rules}. Для PSP или
маркетплейса расчётная модель наследуется от
эквайера, и~менять её в~одностороннем порядке нельзя.
В~российском контуре нетто-позицию для участников ПС~\enquote{Мир} и~для внутрироссийского трафика Visa/Mastercard формирует ОПКЦ~НСПК и~передаёт в~Банк России для расчёта по корсчетам~\cite{nspk-mir-rules,cbr-ps}.

Внутри нетто-свёртки скрыты несколько компонентов: покупки, возвраты,
\textit{interchange fee} (межбанковская комиссия, которую эквайер
платит эмитенту по каждой операции; ставки и~категории
публикуются схемами и~регуляторами~\cite{visa-interchange,eu-ifr}),
\textit{assessments} (сетевые сборы схемы с~эквайера за обработку
объёма) и~штрафные списания. Чтобы сверка работала, каждый компонент
должен быть выделен отдельно. Итоговая сумма по выписке должна
складываться как покупки минус возвраты минус interchange минус
assessments минус прочие удержания. Если не раскрыт хотя бы один компонент,
дельта между ожидаемой и фактической суммой становится
неопознанным расхождением, которое команда расследует вручную каждый
расчётный день.

\begin{figure}[tbp]
  \centering
  \includefiguresvg[width=.7\linewidth]{assets/figures/ch30-reconciliation/net-vs-gross}
  \caption{Нетто- и~валовой расчёт: 100 транзакций в~одну позицию.}\label{fig:net-vs-gross}
\end{figure}

Если схема работает по нетто-модели, искать в~банковской выписке сумму
всех покупок за день бессмысленно. Правильный порядок:
сначала сверять клиринговый поток с~расчётной позицией,
а~потом расчётную позицию с~банковской выпиской.

\practicenote{%
  Не сравнивайте авторизационный журнал сразу с~банковской выпиской:
  авторизация ещё не стала клирингом, выписка уже отражает расчёт.
  Правильный порядок сверки: авторизация против клиринга, клиринг
  против расчётной позиции, расчётная позиция против выписки.

}
\section{Исключения}\label{sec:recon-exceptions}

Транзакции, которые не сходятся автоматически, оставлять в~общей
массе нельзя. Каждое несоответствие нужно классифицировать: понятный тип и~свой маршрут обработки, автоматический или ручной.

Чаще всего первым всплывает расхождение суммы: в~ресторане
авторизовали одну сумму, а~в~клиринге пришла другая, в~гостинице
предварительная блокировка не совпала с~окончательным чеком. Такая
запись требует проверки допустимого отклонения,
и~если порог превышен, запись штатной не считается~\cite{visa-core-rules}.

Другая группа исключений -- клиринг без авторизации. Так бывает при
принудительном представлении, техническом сбое или в~редких сценариях,
когда авторизация прошла вне основного процесса записи. Главное здесь — объяснение происхождения: перед нами либо допустимый сценарий,
либо признак риска.

Обратная проблема -- авторизация без клиринга, когда операция была разрешена,
но в~клиринг так и~не ушла из-за ошибки маршрутизации,
сбоя у~продавца или истечения окна представления. Если не закрыть запись
вовремя, держатель карты увидит только авторизационную
блокировку, а~продавец так и~не получит выручку.

Ещё один сценарий -- дубликаты в~клиринговом файле. Если
эквайер отправляет одну и~ту же запись дважды из-за ошибки
повторной передачи или сбоя батч-процессора, а~сеть не
отфильтровала дубль, эмитент получает два идентичных
требования. Обнаруживает проблему обычно эмитент при сверке,
когда видит два клиринга к~одной авторизации. Стандартный путь
исправления -- возвратный клиринг (\textit{reversal presentment}),
который эквайер отправляет через сеть~\cite{visa-core-rules}.
Для команды сверки дубликат опаснее расхождения в~сумме,
ведь его легко пропустить: каждая запись в~отдельности
выглядит корректной.

Позднее представление (Late Presentment) -- транзакция дошла до схемы
позже применимого срока представления. У
Mastercard это создаёт риск возвратного платежа из-за истечения периода защиты
от спора, а~у Visa до 12 апреля 2024 года было закреплено отдельное
основание спора Condition 12.1 (Late Presentment); с 13 апреля 2024
Visa объединила его с Condition 11.3 \enquote{No Authorization /
Late Presentment}~\cite{mc-tpr,visa-core-rules}. Для команды сверки
одной бухгалтерской классификации недостаточно: нужно
понять, нарушено ли окно представления и не открыло ли это спор по
правилам схемы. В~российском контуре через ОПКЦ~НСПК сроки представления и~основания возвратов закреплены правилами ПС~\enquote{Мир}; коды Visa и~Mastercard здесь не применяются~\cite{nspk-mir-rules}.

\section{DCC и мультивалютный расчёт}\label{sec:recon-dcc}

Устройство DCC, распределение курсового риска и~маршрутизация по коридорам --
в~главе~\ref{ch:multicurrency}. DCC меняет точку фиксации курса, и~часть расхождений рождается из сочетания валюты транзакции, валюты расчёта сети, даты обработки и наценки DCC. Сверка должна уметь отделять курсовое отклонение
от неверной проводки.

После марта 2022~года карты российских эмитентов Visa и~Mastercard за~рубежом не принимаются, поэтому DCC на~стороне держателя российской карты за~пределами РФ неприменима. DCC сохраняет смысл на~стороне российских эквайеров, обслуживающих карты иностранных эмитентов в~ТСП в~России, и~в~кобейджевых сценариях \enquote{Мир--UnionPay} и~\enquote{Мир--JCB}, где конверсию выполняет партнёр-сеть.

Практический список для проверки мультивалютной сверки: (1)~зафиксировать
валюту транзакции (в~Mastercard IPM~--- DE~49), валюту расчёта (DE~50)
и, если применимо, валюту биллинга (DE~51); транзакция с~DCC помечается
DE~54 с~типом суммы~58~\cite{mc-dcc-guide};
(2)~хранить дату и~курс конвертации рядом с~записью; (3)~отделять
наценку DCC от конвертации по правилам схемы; (4)~сверять расчётную позицию
в~валюте расчёта (валюта транзакции для этого шага не подходит). Без любого из четырёх
пунктов система будет регулярно показывать расхождения, которые
объясняются курсовой механикой.

\section{Автоматизация сверки}\label{sec:recon-automation}

Механизм сопоставления находит одну и~ту же
операцию в~разных представлениях. Он работает в~несколько проходов:
сначала строгий матч по идентификаторам и сумме, затем более мягкая
проверка по дате и~другим полям, затем очередь исключений. Чем
слабее критерии сопоставления, тем аккуратнее должна быть ручная
верификация.

Типичная архитектура использует три прохода. Первый~--- строгий:
совпадение по ARN или идентификатору схемы, точная сумма, дата
в~пределах одного дня. Второй -- мягкий: совпадение по сумме и
карте с~допуском по дате в $\pm$2 дня, учёт частичных возвратов.
Третий -- ручная очередь: всё, что не сопоставилось автоматически.
Если на третий проход регулярно приходится больше 2--3\% записей,
проблема не в алгоритме, а~в~качестве входных данных или
в~пропущенном сценарии. Пример строки на каждом проходе с~иллюстративными долями показан на~рис.~\ref{fig:matching-passes}.

У~записи общая ссылка \texttt{reference} -- RRN на стороне авторизации,
ARN или эквивалент в~клиринге; реконсилиатору достаточно того, что строка
одинаково построена в~обоих источниках. Строгий проход ищет точную
пару по \texttt{reference}, сумме и~валюте с~узким временным
допуском; мягкий -- ту же сумму и~валюту с~широким допуском,
\texttt{reference} игнорирует (RRN авторизации и~ARN клиринга могут
не совпасть при позднем представлении) и~выбирает ближайшую
по времени из доступных. Записи, не сопоставленные ни одним проходом,
уходят в \texttt{only\_internal} или \texttt{only\_external} --
это и~есть очередь исключений:

\codefile{Python}{samples/ch30-reconciliation/reconcile.py}

Позднее представление (\textit{late presentment}) -- частая причина ложных
расхождений. Транзакция авторизована вчера, но клиринговая запись
появляется через два-три дня. Если сверка закрывает день жёстко
по календарной дате, такая запись останется сиротой до прихода
клиринга, а~потом потребует ручного сопоставления. Более устойчивый
подход -- оставлять окно ожидания. Несопоставленные авторизации
ждут клиринга в~течение допустимого срока представления схемы
и~только после истечения окна попадают в~очередь исключений.
После реформы Visa Authorization Framework 13~апреля 2024~года эти
сроки таковы: 5~дней для розничного retail card-present и~MIT, 10~дней
для CNP, до~30~дней для T\&E (Travel \& Entertainment -- отели,
аренда авто, круизные линии, авиаперевозки) и~сценариев EAS (Estimated
Authorization -- авторизация на оценочную сумму с~ожидаемой
финализацией позже, помечается полем estimated authorization
indicator)~\cite{visa-core-rules,mc-tpr}.

Три метрики показывают состояние сверки: доля автоматически
сопоставленных записей (match rate), средний возраст
несопоставленной записи и~размер очереди исключений на конец
дня. Конкретные пороги \enquote{здоровых} значений у~каждого
процессора свои и~зависят от стабильности форматов клиринговых
файлов, доли возвратов и качества клиентских реквизитов;
универсальных публичных бенчмарков по этим метрикам в~открытых
источниках нет. Самый надёжный ориентир~--- собственный исторический
ряд. Если match rate падает, причина обычно в~изменении формата
файла или в~новом сценарии, которого парсер не знает. Если
возраст несопоставленных записей растёт, очередь перегружена,
и команда не справляется с~ручным разбором.

\begin{figure}[tbp]
  \centering
  \includefiguresvg[width=.75\linewidth]{assets/figures/ch30-reconciliation/matching-passes}
  \caption{Три прохода сверки: строгий, мягкий, ручная очередь.}\label{fig:matching-passes}
\end{figure}

Готовый сервис сверки имеет смысл там, где нет прямого доступа к
схемам или объём транзакций не оправдывает собственной разработки.
Свой механизм строят там, где есть сложный мультивалютный
процесс, собственная ERP (Enterprise Resource Planning -- учётная
система предприятия) и~специфические правила нетто-расчёта.

\practicenote{%
  Самые упрямые ложные расхождения в~сверке связаны с~датой отсечки.
  Схемный расчёт живёт по операционному времени закрытия дня;
  календарная дата здесь не работает. Поэтому транзакция, прошедшая поздно вечером,
  может попасть в~следующий расчётный день. Если система жёстко
  привязана к~календарной дате, фиктивные расхождения будут появляться
  регулярно, особенно в~конце недели и~перед праздниками.

}

\section{Клиринг и~сверка ПС \enquote{Мир} через НСПК}\label{sec:recon-mir}

Архитектура ОПКЦ (операционный и~платёжный клиринговый центр) НСПК
(Национальная система платёжных карт) и~жизненный цикл внутрироссийской
карточной операции -- в~главе~\ref{ch:mir} (раздел~\ref{sec:mir-opkc}).

Клиринговые файлы НСПК для сверки по картам \enquote{Мир} наследуют ISO~8583 с~собственными расширениями, кодами причин и~идентификаторами~\cite{nspk-mir-rules}.\footnote{Имена файлов, периодичность формирования и~состав полей -- закрытая часть документации НСПК; перед публикацией сверяется с~действующей редакцией интерфейсной спецификации.} Эквайеру
и~эмитенту приходится сверяться сразу на нескольких стыках: между
собственным авторизационным журналом и клиринговой записью НСПК
(типичный случай -- пропущенный или поздно поступивший \textit{reversal}),
между клиринговой позицией ОПКЦ и~расчётным отчётом (комиссии, плата
за услуги ОПКЦ), между расчётным отчётом и~движением по корсчёту
в~Банке России.

Кобейджинговые карты \enquote{Мир--UnionPay} (международная карточная
платёжная система Китая; за пределами Китая представлена дочерней
компанией UnionPay International) и \enquote{Мир--JCB} (Japan Credit
Bureau, японская международная карточная платёжная система) добавляют
параллельный процесс. Транзакции, прошедшие за рубежом
через сеть партнёра, возвращаются в~местный процессинг уже как
зарубежная операция эмитента, со своим клиринговым форматом
UnionPay International или JCB. Эмитенту кобейджинговой программы
приходится держать несколько процессов сверки одновременно.

\section{Сверка платежей по СБП}\label{sec:recon-sbp}

C2B-операции (consumer-to-business -- платежи физического лица в~адрес
бизнеса) по СБП (Система быстрых платежей; см.~гл.~\ref{ch:sbp},
раздел~\ref{sec:sbp-scenarios}) проходят по принципиально другой
схеме, чем карточные. У~СБП нет двухтактной логики \enquote{авторизация
и~потом клиринг}, поскольку исполненный перевод окончателен и~сразу
учтён на корсчёте получателя.
Сверка по СБП поэтому меняет акцент. Она проверяет связку
между платёжной сессией продавца и~фактическим зачислением
на расчётный счёт; сведение авторизационного потока с~клиринговым
здесь не нужно.

Источников данных у~продавца несколько: собственный журнал платёжных
сессий и~заказов; реестр операций от банка-эквайера или PSP,
полученный по протоколу ОПКЦ СБП (раздел~\ref{sec:sbp-protocol});
банковская выписка по расчётному счёту продавца. Здоровая
сверка СБП сопоставляет эти потоки по идентификатору
\texttt{operationId} (или его аналогу в API банка-партнёра),
проверяет совпадение сумм и~по необходимости ловит частичные
возвраты как отдельную цепочку операций.

Финальность перевода означает, что между \enquote{ожидает подтверждения} и \enquote{исполнен} не возникает отдельного клирингового шага. Расхождения, которые
в~карточном мире сглаживаются клиринговым окном, в~СБП проявляются
немедленно и~требуют разбора до закрытия дня. Возврат сопоставляется
с~исходным платежом отдельной операцией, без формы сторно. Сетевого
возвратного платежа в~СБП нет (см.~\ref{sec:disputes-sbp}), поэтому
позиции \enquote{спорная операция} в~реестре сверки не возникает.

B2C-выплаты (business-to-consumer -- выплаты бизнеса в~адрес
физических лиц) по реестру требуют отдельного процесса. Сверка
идёт по реестру в~целом и~по каждой строке отдельно с~проверкой
идемпотентности строк, повторной отправкой только ошибочных
позиций и~отдельным контролем статуса каждого получателя.

\practicenote{%
  Если ваш продукт принимает оплаты и~через карты, и~через
  СБП, держите две независимые системы сверки и~не пытайтесь
  свести их в~одну универсальную модель: общая модель
  обычно теряет важные различия. Общим остаётся только
  формальный объект~--- продажа, с~которой связан один или
  несколько расчётных потоков.
}
